% Demo paper for the meertex website screenshots. Every name, number and
% reference in it is made up.
\documentclass[conference]{IEEEtran}

\usepackage{booktabs}
\usepackage{cite}

\begin{document}

\title{Carbon-Aware Scheduling for Continuous Integration}

\author{\IEEEauthorblockN{Anonymous Authors}
\IEEEauthorblockA{Submission 318}}

\maketitle

\begin{abstract}
Continuous integration runs millions of builds a day, and most of them are not
urgent. We study how much of a project's build energy can move to hours with a
cleaner electricity mix without delaying the developers who wait for the
result. On the build histories of 214 open-source projects, deferring nightly
and documentation jobs by at most six hours cuts their emissions by 31\,\%.
Pull request builds stay where they are.
\end{abstract}

\section{Introduction}\label{sec:intro}

A continuous integration service starts a build for every push. Developers
wait for some of these builds and ignore most of the others. Nightly builds,
dependency updates and documentation jobs can run a few hours later without
anyone noticing. The carbon intensity of the grid changes by a factor of two
over a day in many regions~\cite{brandvold2024grid}.

We ask how much of the build energy of a project could move to cleaner hours,
and what that costs in waiting time.

\section{Method}\label{sec:method}

We replayed the build histories of 214 projects against hourly carbon
intensity data for twelve regions. A job may move when nobody waits for it,
which we approximate by its trigger. Jobs started by a schedule, a bot or a
documentation change may move by up to six hours.

\section{Results}\label{sec:results}

\begin{table}[t]
  \caption{Emissions of movable jobs, before and after deferral.}
  \label{tab:results}
  \centering
  \begin{tabular}{lrr}
    \toprule
    Job type & Share of energy & Saved \\
    \midrule
    Nightly builds & 22\,\% & 34\,\% \\
    Dependency updates & 9\,\% & 29\,\% \\
    Documentation & 4\,\% & 27\,\% \\
    \bottomrule
  \end{tabular}
\end{table}

Table~\ref{tab:results} lists the job types that may move. Together they are
35\,\% of a project's build energy, and deferring them saves 31\,\% of their
emissions.

\begin{thebibliography}{1}
\bibitem{brandvold2024grid}
I.~Brandvold, ``Hourly carbon intensity of twelve electricity grids,'' in
\emph{Proc. Workshop on Sustainable Computer Systems}, 2024, pp. 5--11.
\end{thebibliography}

\end{document}
